% Draft status: V0
% Evidence status: checked where possible; TODOs mark missing evidence.
\section{Problem Definition and Methodology}

Let \(D=\{(x_i,y_i,r_i)\}_{i=1}^{n}\) be an object-recognition dataset, where \(x_i\) is an image, \(y_i\) is the object label, and \(r_i\) is the geographic region associated with the image. Given a model \(f\), this paper evaluates whether object-recognition accuracy is consistent across regions. The unit of analysis is a dataset-model-training-state tuple: one dataset, one model, and either the pretrained or fine-tuned state.

Overall object accuracy is defined as
\[
    \mathrm{Acc}(f)
    =
    \frac{1}{n}\sum_{i=1}^{n}\mathbf{1}\{f(x_i)=y_i\}.
\]
For a region \(r\), per-region accuracy is
\[
    \mathrm{Acc}_{r}(f)
    =
    \frac{\sum_{i=1}^{n}\mathbf{1}\{r_i=r\}\mathbf{1}\{f(x_i)=y_i\}}
    {\sum_{i=1}^{n}\mathbf{1}\{r_i=r\}}.
\]
This allows the same model to be evaluated both as an average recognizer and as a region-conditioned recognizer.

The paper's geographic-disparity metric is the maximum difference between regional accuracies:
\[
    \Delta_{\mathrm{geo}}(f)
    =
    \max_{r}\mathrm{Acc}_{r}(f)
    -
    \min_{r}\mathrm{Acc}_{r}(f).
\]
This metric follows the same intuition as worst-group evaluation: average accuracy can hide the weakest group performance \citep{sagawa2020groupdro}. It also follows the regional-evaluation motivation of prior Dollar Street and GeoDE work \citep{devries2019object,ramaswamy2023geode}. The formula here is this paper's adapted best-minus-worst regional gap; it should not be read as a full fairness metric.

The paper compares two training states. The pretrained state evaluates ImageNet-compatible pretrained models before dataset-specific fine-tuning. The fine-tuned state evaluates the same model family after fine-tuning on the training split for each dataset. The fine-tuning change in average accuracy is
\[
    \Delta_{\mathrm{FT,Acc}}
    =
    \mathrm{Acc}(f_{\mathrm{fine\text{-}tuned}})
    -
    \mathrm{Acc}(f_{\mathrm{pretrained}}).
\]
The fine-tuning change in geographic disparity is
\[
    \Delta_{\mathrm{FT,Disparity}}
    =
    \Delta_{\mathrm{geo}}(f_{\mathrm{fine\text{-}tuned}})
    -
    \Delta_{\mathrm{geo}}(f_{\mathrm{pretrained}}).
\]
A positive \(\Delta_{\mathrm{FT,Acc}}\) means that fine-tuning improves average accuracy. A negative \(\Delta_{\mathrm{FT,Disparity}}\) means that fine-tuning reduces measured regional disparity. A positive \(\Delta_{\mathrm{FT,Disparity}}\) means that fine-tuning increases measured disparity even if average accuracy improves.

The analysis is descriptive. It can show that high average accuracy and low regional disparity do not always coincide. It can also show that fine-tuned and pretrained states differ. It cannot prove why the difference occurs without additional evidence about class balance, region composition, pretraining data, model capacity, and optimization.
